\subsection{Characterization and conjugacy of systems of positive roots}

\begin{lemma}\label{sga3:XXI:3.3.1}
If \(\alpha\in\mathscr{P}(\Delta)\), \(\beta\in\Delta\), and
\(\alpha\) and \(\beta\) are not proportional, then
\[
 s_\beta(\alpha)\in\mathscr{P}(\Delta).
\]
\end{lemma}

\begin{proof}
One has
\[
 s_\beta(\alpha)=\alpha-(\beta^\ast,\alpha)\beta.
\]
At least one simple root other than \(\beta\) occurs with a nonzero,
and hence strictly positive, coefficient in the expression of
\(\alpha\). It occurs with the same coefficient in the expression of
\(s_\beta(\alpha)\); therefore \(s_\beta(\alpha)\) is also positive.
\end{proof}

\begin{corollary}\label{sga3:XXI:3.3.2}
If \(\beta\in\Delta\), then the reflection \(s_\beta\) permutes the
elements of \(\mathscr{P}(\Delta)\) that are not proportional to
\(\beta\).
\end{corollary}

\begin{lemma}\label{sga3:XXI:3.3.3}
Let
\[
 \alpha\in\mathscr{P}(\Delta)\setminus\Delta
\]
be indivisible. There exists \(\beta\in\Delta\) such that
\[
 s_\beta(\alpha)\in\mathscr{P}(\Delta)
 \qquad\text{and}\qquad
 \operatorname{ord}_\Delta\bigl(s_\beta(\alpha)\bigr)
   <\operatorname{ord}_\Delta(\alpha).
\]
\end{lemma}

\begin{proof}
By Lemma~\ref{sga3:XXI:3.1.1}, there is a \(\beta\in\Delta\) such
that \((\beta^\ast,\alpha)>0\). Since \(\alpha\notin\Delta\) and
\(\alpha\) is indivisible, it cannot be proportional to \(\beta\).
Thus \(s_\beta(\alpha)\in\mathscr{P}(\Delta)\) by
Lemma~\ref{sga3:XXI:3.3.1}, and
\[
 \operatorname{ord}_\Delta\bigl(s_\beta(\alpha)\bigr)
 =\operatorname{ord}_\Delta(\alpha)-(\beta^\ast,\alpha)
 <\operatorname{ord}_\Delta(\alpha).
\]
\end{proof}

\begin{corollary}\label{sga3:XXI:3.3.4}
If \(\alpha\in\mathscr{P}(\Delta)\) is indivisible, then there is a
sequence \(\beta_p\in\Delta\), \(p=1,\ldots,q\), such that
\[
 \alpha_p
 =s_{\beta_p}\cdots s_{\beta_1}(\alpha)
 \in\mathscr{P}(\Delta)
 \quad\text{for }p=1,\ldots,q,
\]
and \(\alpha_q\in\Delta\).
\end{corollary}

\begin{proof}
Apply Lemma~\ref{sga3:XXI:3.3.3} repeatedly, inducting on
\(\operatorname{ord}_\Delta(\alpha)\).
\end{proof}

\begin{proposition}\label{sga3:XXI:3.3.5}
The Weyl group is generated by the \(s_\alpha\), for
\(\alpha\in\Delta\). Every indivisible root is conjugate under an
element of the Weyl group to a simple root.
\end{proposition}

\begin{proof}
The second assertion follows immediately from
Corollary~\ref{sga3:XXI:3.3.4}. The first then follows from
Corollary~\ref{sga3:XXI:1.2.10} and
Proposition~\ref{sga3:XXI:2.1.1}.
\end{proof}

\begin{proposition}\label{sga3:XXI:3.3.6}
Let \(R_+\) be a system of positive roots, and let \(P'\subset R\) be
closed and satisfy (P 2). Then there exists \(w\in W\) such that
\[
 w(R_+)\subset P'.
\]
\end{proposition}

We first state the corollaries.

\begin{corollary}\label{sga3:XXI:3.3.7}
The Weyl group acts transitively on the set of systems of positive
roots, and on the set of systems of simple roots.
\end{corollary}

\begin{corollary}\label{sga3:XXI:3.3.8}
A subset of \(R\) is a system of positive roots if and only if it is
closed and satisfies (P 1) and (P 2).
\end{corollary}

\begin{corollary}\label{sga3:XXI:3.3.9}
Suppose that \(\Gamma_0(R)\) is endowed with an ordered-group
structure for which every root is either \(>0\) or \(<0\). Then the
set of positive roots for this order is a system of positive roots.
\end{corollary}
